% !Mode:: "TeX:UTF-8"
%%%%%%%%%%%%%%%%%%%%%%%%%%%%%%%%%%%%%%%%%%%%%%%%%%%%%%%%%%%%%%%%%%%%%%%%%%%%%%% 
%                          _
%  _____ ____ _ _ __  _ __| |___ ___
% / -_) \ / _` | '  \| '_ \ / -_|_-<
% \___/_\_\__,_|_|_|_| .__/_\___/__/
%                    |_|
%  _              _         _   _
% | |__ _  _   __| |_  _ __| |_(_)_ _  __ _  _ ___
% | '_ \ || | / _` | || (_-<  _| | ' \/ _| || (_-<
% |_.__/\_, | \__,_|\_,_/__/\__|_|_||_\__|\_, /__/
%       |__/                              |__/
%%%%%%%%%%%%%%%%%%%%%%%%%%%%%%%%%%%%%%%%%%%%%%%%%%%%%%%%%%%%%%%%%%%%%%%%%%%%%%%



% \vspace{0.8em}
% \section{国内外研究现状及分析}\label{Sec: literature review}
% \vspace{0.8em}


%%%%%%%
\vspace{0.8em}
\section{学位论文的主要研究内容、实施方案及其可行性论证}\label{Sec: content and plan}
\vspace{0.8em}

\vspace{0.8em}
\subsection{主要研究内容}\label{content}
\vspace{0.8em}

基于国内外研究现状的分析，结合深圳市低空经济发展的现实需求，急需对四旋翼无人机的螺旋桨进行优化分析，以提高四旋翼无人机的飞行性能，如减小飞行阻力以实现固定能量供给的条件下的长时间飞行。因此，本文的主要研究内容分为下五点：

（1） 基于简单的理论模型快速生成低保真度的目标螺旋桨的数据集

基于风洞实验的螺旋桨空气动力学参数需要消耗大量的人力物力财力，基于CFD的螺旋桨仿真计算耗时长，单个算例需数天时间，为减少计算时间，本研究计划采用叶素动量理论对螺旋桨进行空气动力学参数计算。叶素动量理论能够在精度不太低的条件下，大幅缩短计算时间，单个螺旋桨算例的计算时间只需数分钟，这样便能快速获得一组用于优化螺旋桨外形的数据集。在优化过程中，螺旋桨的外形需要不断被改变，耗时短的计算方法有着极大的优势，大幅缩短整个优化周期的计算时间，提高优化效率。

%建立鲁棒性的螺旋桨空气动力学参数计算方法的目的是，实现在优化过程中快速计算螺旋桨翼型的空气动力学参数，应用于判断优化循环中新的螺旋桨性能是否提升。在优化的过程中，需要不断调整螺旋桨外形，计算其空气动力学参数，进而测试螺旋桨性能是否得到优化。风洞实验测量需要提供新螺旋桨模型，进而进行测量，需要耗费大量的人力物力财力，不适合联合优化算法计算。采用高保真度的CFD计算同样需要消耗大量的计算资源，每次优化过程涉及到新的几何模型和网格模型构建，工作量巨大，也不适合。


（2）建立基于拓扑特征的数据分析方法

本研究将基于拓扑特征，建立一个分析数据并能够指引如何选择优化算法的方法，该方法可以用来分析数据集的数据结构，也可以用来分析优化过程中的数据，帮助及时发现生成的数据集和优化过程中出现的问题。拓扑特征能够用于分析高维空间数据，揭示数据中隐含的信息，从而帮助分析选择合适的优化算法。现有研究缺少对优化数据的分析方法，优化方法的选择主要基于研究人员的经验，在进行优化过程后，难以快速分析数据的好坏，只能通过实验结果或CFD结果的验证来判断优化结果的好坏。如果在优化过程中出现问题，难以及时发现和改正。

（3）搭建连接低保真性数据与高保真性数据的多重保真性模型

低保真数据获取容易，但是数据的精度低。低保真数据往往来自经验公式的估算、高保真数据精度高，但是获取难度大，主要获取途径有高精度地实验测量、直接数值模拟计算（Direct numerical simulation, DNS）等。本研究计划搭建低保真数据与高保真数据之间的桥梁，建立一个具有多重保真性的模型，既具有低保真数据的获取难度低的有点，又保证数据的高精度性。%多重保真度模型与低保真模型和高保真度模型之间的关系如图\ref{fig:MFM}所示。



（4）建立适用于高维的螺旋桨外形优化算法

本研究计划用于优化螺旋桨外形的输入参数不少于10个，因此该优化过程是一个在高维空间进行寻优的过程。由于高维空间复杂度高，且会引入维数灾难，%~\cite{}%找一个维数灾难的论文，
使高维空间的优化问题较难实现，因此急需一个能够在高维空间中找到最优解的算法。
随着人工智能的发展，这一难题得以初步解决。这些人工智能算法包括人工神经网络（Artificial Neural Network，ANN）、粒子群算法（Particle Swarm Optimization，PSO）、基于“教”与“学”模式的优化算法（teaching-learning-based optimization，TLBO）等。本研究将建立一个基于人工智能的、适用于高维空间的优化算法。该算法将基于拓扑特征的分析结果，进行优化算法的选择、结合与应用。例如，在拓扑特征中分析得知数据集中存在较多极小值时，可以继续使用遗传算法（Genetic Algorithm, GA），而在数据集中存在几个极小值时，则可以尝试使用单纯形法（Downhill Simplex Method, DSM）。


（5）搭建螺旋桨外形优化分析的仿真平台并自动给出螺旋桨外形的优化结果

搭建一个能够自动进行螺旋桨外形优化分析的仿真平台，实现通过多重保真性模型这一桥梁，将输入的快速获取的低保真性数据转化为高保真性数据，
使用适用于高维的螺旋桨外形优化算法，应用基于拓扑特征对优化数据进行分析，以实现对螺旋桨外形进行快速地优化。

仿真平台的输入参数为优化参数，如待优化的几何参数，输出参数为优化目标，如阻力系数$C_d$，升力系数和阻力系数的比值$C_l/C_d$，品质因数（figure of merit, $FM$）等。基于已搭建的螺旋桨外形优化分析的仿真平台，便可进行螺旋桨的外形优化设计。如目标为减小螺旋桨飞行过程中的空气阻力，即可以最小化阻力系数$C_d$为目标进行优化，或以最大化升力系数和阻力系数的比值$C_l/C_d$为目标进行优化，并给出优化后的螺旋桨几何结构。


%（撰写宜使用将来时态，不能只列出论文目录来代替对研究内容的分析论述）
\vspace{0.8em}
\subsection{实施方案及其可行性论证}\label{plan}
\vspace{0.8em}

（1）低保真度螺旋桨空气动力学参数计算方法

为建立具有低保真度的螺旋桨空气动力学参数计算方法，需要对螺旋桨翼型的空气动力学参数进行快速计算，从而计算螺旋桨整体的空气动力学性能。为实现快速计算，应用简化的基于理论的计算模型是比较合适的，本研究计划采用叶素动量理论进行螺旋桨的空气动力学参数计算，该方法应用广泛，经过多年发展，有多种修正方法可以参考应用。在\ref{analysis}节中已经开展基于QBlade软件的初步验证分析，从图\ref{fig:Cd_Cl_60k}至图\ref{fig:Different Re}结果可以看出，QBlade计算螺旋桨翼型的结果与实验结果趋势吻合，能够用于螺旋桨翼型的计算。
为进一步验证QBlade软件能否进行螺旋桨的空气动力学参数计算，本研究采用APC 10×45MR型号螺旋桨，选取在无人机螺旋桨转速一定的条件下，在不同来流风速转速的升力系数、功率系数和螺旋桨效率作为对比参数，将QBlade软件计算结果与Liu Xiran等人\citep{Liu2023AST}的实验数据结果进行验证，对比结果如图\ref{fig:cruise}所示。螺旋桨的转速RPM值和来流风速用进距比归一化，可以看出，QBlade的计算结果与实验结果的趋势相同，可以采用该软件进行低保真度螺旋桨空气动力学参数数据集的生成。


% 在翼型计算方面，本研究计划首先通过调研找到可能的使用方法，再结合课题的实际需求，筛选出合适的方法。在螺旋桨整体计算方面，可以采用叶素动量理论%找一个引用
% 或采用CFD仿真计算\cite{Chauhan2021}。

% 经过前期的调研，现有文献中多采用课题组采用的方法可分为两种，一种为利用实验或CFD计算生成的翼型数据库构建一个代理模型，以计算优化过程中新翼型的空气动力学参数。另一种是直接采用翼型计算程序，如使用课题组长期开发使用的CFD计算程序\cite{Wu2022}或Xfoil。第一种方法建立的代理模型多基于Kriging模型或神经网络模型，训练模型的时间长且难以评估代理模型的准确性。CFD计算程序开发难度大，在课题组未有长期使用的计算程序的前提下独立开发一套计算程序难度大，需要的开发周期长。Xfoil是一个计算二维翼型的空气动力学参数的计算程序，它被广泛应用于翼型或螺旋桨的优化设计研究中\cite{Radi2023, Hoyos2021}。Xfoil基于Fontran编程语言编写，具有计算结果准确\cite{Morgado2016}、计算效率高、易于和其他编程语言建立联合接口\cite{Morgado2016,Radi2023,Seeni2020}的优点。随着内嵌物理知识神经网络（Physics Informed Neural Network, PINN）的发展和Python编程语言的发展，Sharpe\cite{neuralfoil}开发了基于PINN的Python计算库NeuralFoil，用于高效快速的计算翼型的空气动力学参数。本研究计划通过对比Xfoil和NeuralFoil，选出更为高效、易于实现的翼型计算方法。在螺旋桨的整体计算方面，由于本研究计划搭建一个快速的低保真度螺旋桨空气动力学参数的计算方法，所以耗时的CFD仿真计算不在考虑之列。因此本研究计划采用高鲁棒性的叶素动量理论来进行螺旋桨的整体计算。


（2）基于拓扑特征的数据分析方法

数据拓扑分析有利于展现出隐藏在数据背后的逻辑，分析数据点与数据点之间的特征，从而帮助发现数据是否存在噪声等问题，从而快速判断优化数据的可行性。本文计划提出一种基于拓扑特征的数据分析方法，名为持续性数据拓扑（Persistent Data Topology, PDT）。该方法主要分为三个部分。一是极值点的获取，二是拓扑性数据降噪，三是获取最速下降轨迹。目前该方法已经得到多中流动问题的控制优化数据的验证，如螺旋桨外形优化问题、射流的控制和优化问题。

图\ref{fig:propeller}所示为应用PDT分析螺旋桨的几何外形优化的结果图。螺旋桨几何外形优化数据来源于以CFD计算为基准，应用非支配遗传算法（non-dominated sorting genetic algorithm, NSGA-II）进行优化以寻得螺旋桨最佳性能的几何外形\cite{WangTY2023}。该数据选取螺旋桨叶片归一化半径分别为0，0.4，0.7，0.95位置的扭转角作为被优化量，以最大化$FM$值为优化目标。因此螺旋桨的外形优化数据的维度是4，图\ref{fig:propeller}中所示为4维数据降维至2维空间的示意图，降维的方法是经典多维缩放方法（classical multi-dimensional scaling, cMDS），图中标出的红色菱形为极大值，黄色五角星为最大值，蓝色叉号为对应极值点的邻近点。通过PDT的分析发现，极大值在数据的内部，而不在数据空间的边界上，这说明当前的优化算法已经找到了真实的最大值，不需要再向当前数据空间外继续探索。应用弹性响应模型（Elastic Response Model, ERM）对此螺旋桨外形优化数据进行拓扑性数据降噪处理，发现当$\varepsilon$=0.015, 约23\%$FM$值范围时，数据中只有一个最大值，其他5个极大值都被降噪处理掉。



图\ref{fig:jetLES}所示为应用PDT分析控制优化射流仿真计算的结果图。射流的计算是通过大涡模拟（Large Eddy Simulation, LES）进行计算，控制优化方法是贝叶斯优化和深度学习方法相结合的算法，数据维度为22维，优化目标是最小化代价函数。结果显示该数据中共有57个极值点，同时部分通过PDT方法提取的数据极值点有其对应的射流结构含义。这说明，该方法不仅能够成功引用与高维数据的分析，并能够帮助分析、理解对应的物理问题。



图\ref{fig:jetExp}所示为应用PDT分析控制优化射流实验分析的结果图。该射流实验通过12个变形喷嘴和12个微小射流的主动控制频率相位等参数完成优化控制。优化算法为有针对性的位置突变粒子群算法(Particle Swarm Optimization through Targeted, Position-Mutated, Elitism, PSO-TPME)\cite{Shaqarin2023}，数据维度为36维，优化目标是最小化代价函数。为探究该数据的结构，采用将数据旋转12次并做一次对称变化的方式增强数据，并采用POD进行数据的降维处理，从图\ref{fig:jetExp}可以看出，增强后的数据具有高度对称性，原数据中的最小值在增强变换后变为极小值。




通过上述三组数据的验证，PDT能够成功进行数据的拓扑分析，即提取数据中的极值点，必要时进行拓扑性数据降噪，分析极大值点或极小值点的数据相关性。在高维数据的处理方面，目前PDT能够处理高达36维的数据。因此应用该方法进行本研究的拓扑数据处理是可行的。


（3）高维空间的优化算法

Chernukhin\cite{Oleg2013}指出，螺旋桨翼型的气动形状优化设计问题是一个具有多峰值特征、多模态的优化问题。因而螺旋桨外形的优化设计逐渐朝着优化数据维数越来越高的方向发展，这主要是因为设计空间维数的增加可以使得计算精度提升，这同时也得利于计算机性能的增强。设计参数的增加虽然可以增加更多可能的螺旋桨翼型，但是同时也会使优化算法陷入局部最优\cite{Nikhil2013}。

本研究的优化问题是一个复杂的高维优化问题，选择哪种优化算法是比较困难的。通过拓扑分析可以明确在高维数据中的极值点个数，从而帮助选择适合的优化算法。如果数据中只有一个极值点，则可以采用梯度下降法即可快速找到最优解，如果数据中有几个极值点，则可以采用随机梯度下降法，如果数据中有多于十个的极值点，则可以采用如遗传算法、粒子群算法等的智能优化算法，这些算法已被广泛应用并取得较好的优化结果。
% 在本文计划采用非梯度的单纯形算法，并对单纯形退化现象进行修正，修正方法示意图如\ref{fig:rDSM}所示。提出一种鲁棒性的单纯形优化算法(Robust Downhill Simplex Method, rDSM)，处理高维数据的优化问题。单纯形方法最早由Nelder于1965年提出\cite{Nelder1965}。


%%%%%%%%%%%


\vspace{0.8em}
\section{论文进度安排，预期达到的目标}\label{Sec: schedule}
\vspace{0.8em}


\vspace{0.8em}
\subsection{进度安排}\label{schedule}
\vspace{0.8em}

2023.03 - 2023.6：学习拓扑分析相关知识，初步形成优化数据拓扑分析方法的构建。

2023.07 - 2023.11：完成优化数据拓扑分析方法的提出并使用课题组内同学的优化数据进行应用。调研螺旋桨外形优化相关论文，确定低保真度数据生成方法、多重保真度模型搭建方法、螺旋桨外形优化方法，并进行初步尝试，撰写开题报告国内外研究现状部分。

2023.12 - 2024.02：学习使用QBlade软件并于论文中数据结果进行对比验证。完成开题报告的撰写主要研究内容、研究方案及前期工作结果部分。

2024.03 - 2024.05：使用QBlade软件生成大量螺旋桨的低保真度数据。完成开题报告撰写，完成开题答辩。

2024.06 - 2024.08：应用数据的拓扑分析方法分析低保真数据，根据拓扑分析的结果建立适用于高维计算的优化算法。

2024.09 - 2024.11：用低保真度数据对优化算法进行尝试和验证以及分析。并为获取高保真度数据做准备。

2024.12 - 2025.02：获取高保真度数据，并搭建多重保真度模型。

2025.03 - 2025.05：应用多重保真度模型从低保真度数据获取更为精确的数据，并用该数据进行优化分析。

2025.06 - 2025.08：整合数据拓扑分析的代码、优化算法的代码、多重保真度数据的代码，形成一个完整的螺旋桨外形优化分析的仿真平台，实现输入待优化螺旋桨参数和优化目标，自动输出优化结果的功能。

2025.09 - 2025.11：应用螺旋桨外形优化分析的仿真平台进行螺旋桨在不同数据维数、不同的优化目标条件下的优化分析，并整理结果。

2025.12 - 2026.02：完成博士毕业论文的撰写。

2026.03 - 2026.06：完成博士毕业预答辩和最终答辩。


\vspace{0.8em}
\subsection{预期达到的目标}\label{objectives}
\vspace{0.8em}

完成博士毕业论文，形成2-4篇高水平SCI论文。

%%%%%%%%%
\vspace{0.8em}
\section{学位论文预期创新点}\label{Sec: innovations}
\vspace{0.8em}

（1） 首次在数据处理过程中考虑拓扑结构分析，尤其是应用于高维数据的拓扑特征分析，从中提取到数据中的极值点信息，并用于指导优化算法的选择。

（2）搭建联系低保真度的螺旋桨空气动力参数数据与高保真度的螺旋桨空气动力参数数据之间的多重保真度模型，实现从低保真度数据快速参数分布趋势，通过多重保真度模型获取更为精确的数据，缩短获取高保真数据的时间。
    
    %（3） 提出鲁棒性的单纯形法（rDSM）用来进行高维数据的优化分析。 
    
（3） 搭建一个具有鲁棒性的无人机螺旋桨优化分析的仿真平台，该平台能够实现输入待优化的螺旋桨几何参数以及优化目标侯，自动获取优化结果。


%%%%%%%%%%
\vspace{0.8em}
\section{为完成课题已具备和所需的条件、外协计划及经费}\label{Sec: support}
\vspace{0.8em}

指导老师Noack教授在人工智能、机器学习领域有着丰富的经验，在优化算法的开发和应用方面取得一定的成绩\cite{LiA2022jfm,Shaqarin_PSOE}，足以指导螺旋桨外形优化算法的相关工作。同时，Noack教授在拓扑结构问题方面也取得一定的成绩，可以对数据的拓扑特征方面进行指导\cite{Fernex2020prf}。来自西北工业大学的邬晓敬老师和来自华南理工大学的杨延年老师都是与Noack教授的合作者。他们在无人机领域有丰富的经验，能够在无人机螺旋桨计算方面给予支持。与此同时，Noack教授还与QBlade开发者David Martin有合作，能够帮助解决使用该软件中遇到的问题。此外，课题组提供高性能工作站可以用来处理大数据量处理与计算工作。

课题组经费充足，包括国家自然科学基金（Grants No. 
%12172111, 
12172109和12302293）,广东省基础与应用基础研究基金（ No. 2022A1515011492）以及深圳市科学技术基金（No. JCYJ20220531095605012）.

%%%%%%%%%%%%%
\vspace{0.8em}
\section{预计研究过程中可能遇到的困难、问题，以及解决的途径}\label{Sec: problems}
\vspace{0.8em}

在搭建螺旋桨外形优化分析的仿真平台搭建过程中，可能涉及多种计算程序或平台的接口搭建问题。该问题可以通过创建相对应的接口链接进行数据传递，如果涉及到自行开发的程序设计及编写问题时，应只使用一种编程语言，以方便进行数据传递与计算。

在建立多重保真度模型的过程中，可能因数据量巨大、低保真度数据和高保真数据之间单位等信息不匹配的问题。数据量大、数据维数高导致计算速度下降的问题可使用高性能工作站、超算等方法解决。在建立低保真度数据和高保真度数据时应尽早考虑数据之间的联系，也可以参照高保真数据集的建立方式建立相对应的低保真度数据。
% 在进行优化算法的建立过程中，难点在于高维问题的寻优过程的设计。本研究计划通过使用测试函数在不同维条件下进行验证，由低维测试函数到高维测试函数进行测试。比如本研究计划提出的rDSM优化算法已通过10维优化算法的测试验证。

在进行基于拓扑特征的数据分析时，需要利用数据点的多个邻近点建立凸包，高维数据的凸包建立是一个难点。被广泛应用的qhull算法可以建立10维凸包，但是如果数据的维数增加到20甚至更大，那么qhull不再适用。可以通过拓扑数据分析的相关文献以及代码寻找高维凸包的建立方法，如All Vertex Triangle Algorithm(AVTA)算法\cite{Awasthi2020}。


%%%%%%%%%%%%
\vspace{0.8em}
\section{主要参考文献}\label{Sec: reference}
\vspace{0.8em}

\bibliographystyle{hithesis}
\bibliography{reference}

% Local Variables:
% TeX-master: "../report"
% TeX-engine: xetex
% End: